\chapter{The Politics of Water: A Lifeline for an Island Nation}

\section{A City Without Rivers}

When I first arrived in Singapore, one thing surprised me: this city has almost no natural water resources.

For a tropical island with abundant rainfall, this sounds contradictory. But in fact, Singapore's water problem is longstanding and fundamental --- the island has no large natural rivers, rainwater is difficult to store, and at the time of independence, almost all freshwater supply depended on imports from neighboring Malaysia.

Water is Singapore's most original and most enduring strategic anxiety.

\section{Reservoirs: Decades of Effort After Independence}

During my stay in Singapore, I made a point one morning of visiting a reservoir in the northwestern part of the island. It was around six in the morning, and the sky had not yet fully brightened, but already quite a few residents were exercising along the waterfront --- some walking, some jogging, moving at a leisurely pace.

This body of water is considered quite large by Singapore's standards. The water was clear, and purpose-built walkways lined the banks; from a distance they looked like wooden boardwalks, but up close they turned out to be plastic composite material --- durable and weather-resistant, just like this country's pragmatic approach to everything it does.

The scenery was genuinely pleasant, but behind this tranquil surface lay decades of extraordinarily difficult effort.

In the early years after independence, Singapore had almost no large reservoirs. At that time, large quantities of rainwater simply flowed directly into the sea after falling, wasted. After independence, the Singapore government spent decades systematically planning and building a network of reservoirs to collect scattered rainfall and accumulate it over time. The technical and financial investment in this effort far exceeds what the average person might imagine.

\section{Three Sources, One Goal}

Today, Singapore's water supply comes mainly from three channels.

\textbf{First, purchasing water from Malaysia.} This is the historically dominant source and Singapore's greatest long-term strategic vulnerability. The water supply agreements between the two countries have been repeatedly troubled, and price negotiations have more than once triggered diplomatic friction. For this reason, reducing dependence on Malaysia's water supply has always been one of the core policy goals of every Singapore government.

\textbf{Second, rainwater collection.} Singapore uses a catchment system spread across the entire island to channel rainfall into various reservoirs. Marina Reservoir is perhaps the most representative engineering achievement --- it collects and purifies stormwater runoff from the central urban area, transforming what was once a bay that drained directly into the sea into a freshwater storage facility.

\textbf{Third, seawater desalination.} Singapore has built several desalination plants that use reverse osmosis technology to convert seawater into potable water. This approach is energy-intensive and costly, but for a city-state surrounded by sea on all sides, it has an irreplaceable strategic value.

In addition, Singapore has invested heavily in wastewater recycling technology, injecting highly treated reclaimed water (NEWater) into reservoirs or using it directly for industrial purposes --- incorporating ``waste water'' back into the water cycle and making full use of every drop.

\section{Water Use: Thrift as a Culture}

In Singapore, conserving water is not only a government policy --- it is a way of life deeply embedded in daily routines. The government has long used tiered water pricing, public education, and technical monitoring to build a society-wide awareness of water conservation. Even in years of abundant rainfall, the tone of thrift never changes --- because behind every drop of rain, someone remembers how scarce water once was on this island.

\section{Water and Sovereignty}

Water resources in Singapore are never merely a technical matter. They are deeply intertwined with national sovereignty, diplomatic relations, and national identity.

During his lifetime, Lee Kuan Yew stated on multiple occasions that if Malaysia were to cut off water to Singapore, Singapore would treat it as an act of war. This was not empty rhetoric --- it was a small country with a clear-eyed awareness of its own vulnerability, expressing its existential bottom line.

For this reason, Singapore's sustained investment in water resources over decades is less a matter of infrastructure development than a strategic act of self-preservation for survival. Every step taken toward water self-sufficiency is an effort by this nation to reduce its dependence on the outside world and consolidate its own foundations for existence.

On a resource-scarce island, water is the finest edge of the blade.
